% ============================================================
%  Lernzettel Template (my_dhbw Stil, analysis.tex)
%  Eine Spalte, mit TOC, accent-Theme.
%  Renderer: pdflatex (zweimal)
% ============================================================
\documentclass[11pt, a4paper]{article}

\usepackage[ngerman]{babel}
\usepackage{fontspec}
\setmainfont[
  Path=/usr/share/fonts/TTF/,
  Extension=.ttf,
  BoldFont=DejaVuSans-Bold,
  ItalicFont=DejaVuSans-Oblique,
  BoldItalicFont=DejaVuSans-BoldOblique
]{DejaVuSans}
\setsansfont[
  Path=/usr/share/fonts/TTF/,
  Extension=.ttf,
  BoldFont=DejaVuSans-Bold,
  ItalicFont=DejaVuSans-Oblique,
  BoldItalicFont=DejaVuSans-BoldOblique
]{DejaVuSans}
\setmonofont[Path=/usr/share/fonts/TTF/,Extension=.ttf]{DejaVuSansMono}
\usepackage[top=2cm, bottom=2cm, left=2cm, right=2cm, footskip=1cm]{geometry}
\usepackage{amsmath, amssymb}
\usepackage{xcolor}
\usepackage{titlesec}
\usepackage{titletoc}
\usepackage{enumitem}
\usepackage{fancyhdr}
\usepackage{lastpage}
\usepackage{microtype}
\usepackage{parskip}

\usepackage{booktabs}
\usepackage{array}
\usepackage{longtable}
\usepackage{tabularx}
\usepackage{ltablex}
\keepXColumns
\usepackage{colortbl}
\usepackage[most,breakable]{tcolorbox}
\usepackage{listings}
\usepackage[normalem]{ulem}
\usepackage{pdflscape}
\usepackage[colorlinks=true, linkcolor=accent, urlcolor=accent]{hyperref}

\renewcommand{\familydefault}{\sfdefault}

% ── Theme (my_dhbw accent) ────────────────────────────────────
\definecolor{accent}{RGB}{0, 60, 113}
\definecolor{primaryblue}{RGB}{0, 60, 113}
\definecolor{codebg}{RGB}{245, 245, 245}
\definecolor{figurebg}{RGB}{232, 241, 255}
\definecolor{tablehintbg}{RGB}{240, 255, 240}
\definecolor{solutionbg}{RGB}{242, 255, 242}
\definecolor{rulegray}{RGB}{180, 180, 180}
\definecolor{tableheadbg}{RGB}{0, 60, 113}

% ── Header/Footer ─────────────────────────────────────────────
\pagestyle{fancy}
\fancyhf{}
\renewcommand{\headrulewidth}{0pt}
\renewcommand{\footrulewidth}{0.4pt}
\renewcommand{\sectionmark}[1]{\markboth{#1}{}}
\fancyfoot[L]{\footnotesize\color{accent}\nouppercase{\leftmark}}
\fancyfoot[R]{\footnotesize\color{accent}\thepage\,/\,\pageref{LastPage}}
\fancypagestyle{plain}{%
  \fancyhf{}%
  \renewcommand{\headrulewidth}{0pt}%
  \renewcommand{\footrulewidth}{0.4pt}%
  \fancyfoot[L]{\footnotesize\color{accent}\nouppercase{\leftmark}}%
  \fancyfoot[R]{\footnotesize\color{accent}\thepage\,/\,\pageref{LastPage}}%
}

% ── Typografie ────────────────────────────────────────────────
\titleformat{\section}
    {\color{accent}\large\bfseries}{}{0pt}{}
    [\vspace{-4pt}\color{accent}\rule{\linewidth}{2pt}]
\titleformat{\subsection}
    {\normalsize\bfseries\color{accent!85!black}}{}{0pt}{}
    [\vspace{-3pt}\color{accent!40}\rule{\linewidth}{0.4pt}]
\titleformat{\subsubsection}
    {\small\bfseries\color{accent!70!black}}{}{0pt}{}{}
\titlespacing*{\section}{0pt}{12pt}{4pt}
\titlespacing*{\subsection}{0pt}{8pt}{2pt}
\titlespacing*{\subsubsection}{0pt}{6pt}{1pt}

\setlength{\parindent}{0pt}
\setlength{\parskip}{4pt}
\renewcommand{\arraystretch}{1.25}
\setlist[itemize]{noitemsep, topsep=3pt, parsep=0pt, leftmargin=1.4em}
\setlist[enumerate]{noitemsep, topsep=3pt, parsep=0pt, leftmargin=1.6em}
\setlist[itemize,2]{label=\textendash, leftmargin=1.4em}

% ── Boxen ──────────────────────────────────────────────────────
\newtcolorbox{figurebox}[1]{
  colback=figurebg, colframe=accent!50,
  title={Abbildung: #1}, fonttitle=\small\bfseries,
  breakable, before skip=6pt, after skip=6pt
}
\newtcolorbox{tablehintbox}[1]{
  colback=tablehintbg, colframe=green!50!black!50,
  title={Tabelle: #1}, fonttitle=\small\bfseries,
  breakable, before skip=6pt, after skip=6pt
}
\newtcolorbox{solutionbox}[1]{
  colback=solutionbg, colframe=green!60!black!60,
  title={#1}, fonttitle=\small\bfseries,
  breakable, before skip=6pt, after skip=6pt
}

\lstset{
  basicstyle=\ttfamily\small,
  backgroundcolor=\color{codebg},
  breaklines=true,
  frame=single, framerule=0pt,
  xleftmargin=4pt, xrightmargin=4pt,
  aboveskip=6pt, belowskip=6pt,
  keepspaces=true
}

% ── TOC-Look ──────────────────────────────────────────────────
\renewcommand{\contentsname}{\color{accent}Inhalt}

\providecommand{\nichttitle}{Lernzettel}
\providecommand{\nichtsubtitle}{}

\renewcommand{\nichttitle}{Numerische Methoden}
\renewcommand{\nichtsubtitle}{Lernzettel}


\begin{document}

\begin{center}
{\bfseries\Large\color{accent}\nichttitle}\\[2pt]
\color{accent}\rule{0.35\linewidth}{1pt}\color{black}
\end{center}
\vspace{2pt}

{\small\tableofcontents}
\newpage

% ── BODY ──────────────────────────────────────────────────────

\section{Numerische Methoden — Gesamtzettel}


\subsection{1. Grundlagen}

\textbf{Numerische Methoden}: Verfahren zur näherungsweisen Lösung mathematischer Probleme, die analytisch nicht oder nur unpraktikabel lösbar sind.

\begin{itemize}
  \item f(·) = exakte Lösung; f̃(·) = numerische Näherung
  \item x = exakter Input; x̃ = gestörter Input
  \item Notwendig bei: nichtlinearen Gleichungen, transzendenten Funktionen, großen LGS
\end{itemize}

\textbf{Transzendente Funktion}: Funktion ohne endliche Kombination algebraischer Operationen — keine geschlossene Lösung.

\begin{itemize}
  \item Beispiel Sigmoid: σ(x) = 1/(1+e⁻ˣ)
\end{itemize}

\textbf{Smooth}: Funktion mit Ableitungen aller Ordnungen.


\textbf{Moore's Law}: Rechenleistung verdoppelt sich ca. alle 2 Jahre.


\textbf{[a,b]}: geschlossenes Intervall (inkl. Endpunkte).



\subsection{2. Diskretisierung}

\textbf{Diskretisierung}: Zerlegung kontinuierlicher Domänen (Zeit, Raum) in endliche Schritte/Gitter.

\begin{itemize}
  \item Kontinuierlich: f(x), x ∈ [a,b]
  \item Diskret: f(xᵢ), xᵢ = a + ih, i = 0,…,N
  \item Schrittweite: h = (b−a)/N
\end{itemize}


\subsection{3. Floating-Point Arithmetik}


\subsubsection{3.1 Fehlertypen}


{\small
\begin{tabularx}{\linewidth}{XX}
\hline
\rowcolor{tableheadbg}
{\color{white}\textbf{Fehler}} & {\color{white}\textbf{Ursache}} \\[2pt]
\hline
\endhead
\hline
\endfoot
Modeling Error & Modell ist Vereinfachung der Realität \\
Truncation/Approximation & Endliche Approximation unendlicher Prozesse \\
Rounding Error & Endliche Stellenzahl bei Speicherung \\
Overflow & Magnitude zu groß → ±∞ \\
Underflow & Magnitude nahe Null → 0 \\
\end{tabularx}
}



\subsubsection{3.2 Floating-Point Darstellung}

\textbf{Formel}: x = (−1)ˢ · m · 2ᵉ — Standard IEEE 754.

\begin{itemize}
  \item s = Vorzeichenbit, m = Mantisse (Präzision), e = Exponent (Skalierung)
\end{itemize}


{\small
\begin{tabularx}{\linewidth}{XXXXX}
\hline
\rowcolor{tableheadbg}
{\color{white}\textbf{Format}} & {\color{white}\textbf{Bits}} & {\color{white}\textbf{Sign}} & {\color{white}\textbf{Exponent}} & {\color{white}\textbf{Mantisse}} \\[2pt]
\hline
\endhead
\hline
\endfoot
HALF & 16 & 1 & 6 & 9 \\
SINGLE & 32 & 1 & 8 & 23 \\
DOUBLE & 64 & 1 & 11 & 52 \\
\end{tabularx}
}


\begin{itemize}
  \item IEEE 754 Single: Bias = 127, e = E − 127; Exponentbereich 0…255, zentriert um 0
  \item Größte Single ≈ 10³⁸ (log₁₀(2¹²⁸) ≈ 38.5)
\end{itemize}

\textbf{Mantisse}: t Bits → 2ᵗ Werte zwischen zwei Zweierpotenzen.

\begin{itemize}
  \item 2-Bit Beispiel: 0.00=0, 0.01=0.25, 0.10=0.5, 0.11=0.75
\end{itemize}

\textbf{Machine Epsilon (ε\_mach)}: Kleinste positive Zahl mit 1 + ε\_mach ≠ 1.

\begin{itemize}
  \item ε\_mach = 2⁻ᵗ
  \item Single: ≈ 1.19 × 10⁻⁷; Double: ≈ 2.22 × 10⁻¹⁶
  \item Absolute Präzision: ε\_mach · |x| (große Zahlen → größere Lücken)
\end{itemize}

\textbf{Fixed-Point}: Vorab feste Bits für Integer- und Fraktionsteil — konstante Präzision unabhängig von Magnitude.

\begin{itemize}
  \item Beispiel: Skalierung 10⁶, 1.234567 → 1234567; max. Rundungsfehler 0.5 × 10⁻⁶
\end{itemize}


\subsubsection{3.3 Loss of Significance}

\textbf{Catastrophic Cancellation}: Rundungsfehlerverstärkung durch Subtraktion fast gleicher Zahlen — führende Stellen löschen sich, nur unzuverlässige niedrigwertige Stellen bleiben.

\begin{itemize}
  \item Beispiel 8 sig. Digits: a=12345678.5, b=12345678.0 → ã−b̃ = 1 statt 0.5
  \item f(ε) = 1/(1+ε−1): bei kleinem ε divergiert Resultat
\end{itemize}

\textbf{Pseudo-Accuracy}: Künstlich hohes Detaillevel, das ungerechtfertigte Genauigkeit suggeriert.



\subsection{4. Fehleranalyse — Charakterisierung numerischer Methoden}


{\small
\begin{tabularx}{\linewidth}{XXX}
\hline
\rowcolor{tableheadbg}
{\color{white}\textbf{Begriff}} & {\color{white}\textbf{Formel}} & {\color{white}\textbf{Beschreibung}} \\[2pt]
\hline
\endhead
\hline
\endfoot
\textbf{Conditioning} κ & κ = \textbackslash{} & f(x) − f(x̃)\textbackslash{} \\
\textbf{Stability} & \textbackslash{} & f̃(x̃) − f̃(x)\textbackslash{} \\
\textbf{Consistency} & \textbackslash{} & f̃(x) − f(x)\textbackslash{} \\
\textbf{Convergence} & \textbackslash{} & f̃(x̃) − f(x)\textbackslash{} \\
\end{tabularx}
}


\begin{itemize}
  \item κ unabhängig von numerischer Methode
  \item Konvergenz strebt meist gegen 0; Nicht-Null = systematischer Bias
  \item Stability = Conditioning des Systems für h → 0
\end{itemize}


\subsection{5. Brute-Force / Grid Search}

\textbf{Grid Search}: Funktion an allen Gitterpunkten auswerten, bestes Ergebnis wählen.

\begin{itemize}
  \item Komplexität: O(Nᵈ)
  \item Extrema auf [a,b] können an kritischen Punkten ODER Endpunkten liegen — beide prüfen
\end{itemize}


\subsubsection{5.1 Beispiel: min sin(x) auf [−3,1]}
\begin{itemize}
  \item Analytisch: f'(x) = cos(x) = 0 ⟹ x* = −π/2 ≈ −1.5708, sin(−π/2) = −1
  \item Grid h=1: sin(−3)≈−0.1411, sin(−2)≈−0.9093, sin(−1)≈−0.8415, sin(0)=0, sin(1)≈0.8415
  \item Numerisches Min: −0.9093 bei x=−2; Approximationsfehler: 0.0907
\end{itemize}


\subsubsection{5.2 LGS-Beispiel}
[2,1;5,1]·[w,b]ᵀ = [11,13]ᵀ → w=2/3, b=29/3

\begin{itemize}
  \item Grid w,b ∈ \{0,2.5,5,7.5,10\}; error₁=|2w+b−11|, error₂=|5w+b−13|
  \item Min bei w=0, b=10, akkumulierter Fehler = 4
  \item Substitution skaliert nicht für große Systeme
\end{itemize}


\subsubsection{5.3 Conditioning/Stability auf min sin(x), x̃ = x ± 0.25}
\begin{itemize}
  \item Conditioning: κ(−1, +0.25) ≈ 0.1599; κ(−1, −0.25) ≈ 0.1074 (gut); κ(0, ±0.25) ≈ 0.2474 (schlecht)
  \item Stability h=1: s ≈ 0; h=0.2 bei x=−1: s ≈ 0.4964
  \item Consistency: c₁ ≥ 0.15853; c₀.₂ ≥ 0.05102 (besser bei kleinerem h)
  \item Beide h konvergieren gleich; h=0.2 verbessert nicht wegen Stabilitätsproblemen
\end{itemize}


\subsection{6. Taylor-Expansion}

\[
f(x) = \sum_{k=0}^{\infty} \frac{f^{(k)}(x_n)}{k!}(x-x_n)^k
\]



{\small
\begin{tabularx}{\linewidth}{XXX}
\hline
\rowcolor{tableheadbg}
{\color{white}\textbf{Ordnung}} & {\color{white}\textbf{Polynom P(x)}} & {\color{white}\textbf{Bedingung}} \\[2pt]
\hline
\endhead
\hline
\endfoot
0. & f(xₙ) & P(xₙ) = f(xₙ) \\
1. & f(xₙ) + f'(xₙ)(x−xₙ) & P'(xₙ) = f'(xₙ) \\
2. & + f''(xₙ)/2·(x−xₙ)² & P''(xₙ) = f''(xₙ) \\
\end{tabularx}
}


\begin{itemize}
  \item Division durch k! kompensiert Mehrfachableitung
  \item \textbf{Lokale} Approximation, verankert an xₙ
\end{itemize}


\subsection{7. Newton-Raphson Methode}

\textbf{Ziel}: x\textit{ finden mit f(x}) = 0 mittels lokaler Steigungsinformation.


\textbf{Lineare Approximation} (1. Ordnung Taylor): f(x) ≈ f(xₙ) + f'(xₙ)·(x − xₙ)


\textbf{Iteration}:

\[
x_{n+1} = x_n - \frac{f(x_n)}{f'(x_n)}
\]


\begin{quote}
Linear genügt, weil wir f(x*)=0 suchen, nicht f(x) approximieren — erster nicht-konstanter Term mit Richtungsinformation.
\end{quote}


\subsubsection{7.1 Algorithmus 1 (Newton-Raphson)}
\textbf{Require}: x₀, f, f', Toleranz ε

\begin{enumerate}
  \item x ← x₀
  \item while |f(x)| > ε:
  \item d ← f'(x)
  \item if |d| < ε\_machine: error "Derivative too small"
  \item x ← x − f(x)/d
  \item return x
\end{enumerate}


\subsubsection{7.2 Quadratische Konvergenz}
\begin{itemize}
  \item eₙ = xₙ − x*
  \item 2. Ordnung Taylor um x\textit{: f(xₙ) = f'(x})·eₙ + f''(ξ)/2·eₙ²
  \item $e_{n+1} \approx -\frac{f''(\xi)}{2f'(x^*)} e_n^2$
  \item Allgemein: |eₙ₊₁| ≤ C·|eₙ|²
  \item Fehler quadriert pro Iteration: 10⁻¹ → 10⁻² → 10⁻³ → 10⁻⁶ → 10⁻¹²
\end{itemize}


\subsubsection{7.3 Beispiel √2, x₀ = 2}


{\small
\begin{tabularx}{\linewidth}{XXXXX}
\hline
\rowcolor{tableheadbg}
{\color{white}\textbf{n}} & {\color{white}\textbf{xₙ}} & {\color{white}\textbf{\textbackslash{}}} & {\color{white}\textbf{eₙ\textbackslash{}}} & {\color{white}\textbf{}} \\[2pt]
\hline
\endhead
\hline
\endfoot
0 & 2.000000000 & 5.86 × 10⁻¹ &  &  \\
1 & 1.500000000 & 8.58 × 10⁻² &  &  \\
2 & 1.416666667 & 2.45 × 10⁻³ &  &  \\
3 & 1.414215686 & 2.13 × 10⁻⁶ &  &  \\
4 & 1.414213562 & 4.5 × 10⁻¹² &  &  \\
\end{tabularx}
}



\subsubsection{7.4 Beispiel sin(x)=0, x₀=3.5}
\begin{itemize}
  \item x₁ = 3.5 − (−0.351)/(−0.936) = 3.125
  \item x₂ ≈ 3.14159; nach 2 Iterationen Fehler < 10⁻⁵
  \item Vgl. Grid h=0.3: 6 Auswertungen, Fehler ≈ 0.042
\end{itemize}


\subsubsection{7.5 Beispiel f(x)=x³−x, x₀=1.4 (f'=3x²−1)}
\begin{itemize}
  \item x₁ = 1.4 − 1.344/4.88 ≈ 1.127
  \item x₂ ≈ 0.976
  \item x₃ ≈ 1.014
\end{itemize}


\subsubsection{7.6 Failure Modes}
\begin{itemize}
  \item \textbf{Zero derivative}: f'(xₙ)=0 → horizontale Tangente, kein Schnittpunkt
  \item \textbf{Oszillation}: bei symmetrischen Funktionen Zyklus
  \item \textbf{Ambiguous starting point}: x₀ bestimmt, welche Wurzel gefunden wird
  \item \textbf{Lokales Optimum}: konvergiert zu nahem, nicht globalem Optimum
  \item \textbf{Divergenz}: bei Ableitung nahe Null
\end{itemize}


\subsection{8. Secant Methode}

\textbf{Motivation}: f'(x) nicht verfügbar (teuer, Black-Box, nicht differenzierbar).


\textbf{Ableitungs-Approximation}: $f'(x_n) \approx \frac{f(x_n) - f(x_{n-1})}{x_n - x_{n-1}}$


\textbf{Iteration}:

\[
x_{n+1} = x_n - f(x_n) \cdot \frac{x_n - x_{n-1}}{f(x_n) - f(x_{n-1})}
\]


\begin{itemize}
  \item Benötigt \textbf{zwei} Startwerte x₀, x₁
  \item Geometrisch: Sekantensteigung statt Tangente
\end{itemize}


\subsubsection{8.1 Algorithmus 2 (Secant Method)}
\textbf{Require}: x₀, x₁, f, ε

\begin{enumerate}
  \item while |f(x₁)| > ε:
  \item s ← (f(x₁) − f(x₀))/(x₁ − x₀)
  \item x\_old ← x₁
  \item x₁ ← x₁ − f(x₁)/s
  \item x₀ ← x\_old
  \item return x₁
\end{enumerate}


\subsubsection{8.2 Beispiel f(x)=x³−x, x₀=1.2, x₁=1.4}
\begin{itemize}
  \item x₂ ≈ 1.4 − 0.744·(0.2/0.616) ≈ 1.159
  \item x₃ ≈ 1.011
\end{itemize}


\subsection{9. Vergleich der Methoden}


{\small
\begin{tabularx}{\linewidth}{XXXXX}
\hline
\rowcolor{tableheadbg}
{\color{white}\textbf{Methode}} & {\color{white}\textbf{Info-Bedarf}} & {\color{white}\textbf{Startwerte}} & {\color{white}\textbf{Konvergenz}} & {\color{white}\textbf{Komplexität}} \\[2pt]
\hline
\endhead
\hline
\endfoot
Grid Search & nur f(x) & Intervall & linear/blind & O(Nᵈ) \\
Newton-Raphson & f(x), f'(x) & 1 (x₀) & quadratisch & wenige Iter., schneller \\
Secant & nur f(x) & 2 (x₀, x₁) & superlinear & kein f' nötig \\
\end{tabularx}
}



\subsection{10. Anwendung — ML/DL Bezug}

\begin{itemize}
  \item Newton 1. Ordnung = \textbf{Gradient Descent} in ML/DL
  \item \textbf{Backpropagation}: verteilt Gradient durch Netz-Schichten
  \item Lokale Methoden finden lokale Lösungen — Konvergenz zum nächsten Optimum, kann bei f'=0 oszillieren/divergieren
\end{itemize}




% ──────────────────────────────────────────────────────────────

\end{document}
